\section{Every proper relative canonical coefficient belongs to the original regular ring}
\label{sec:proper-canonical-coefficient-regularity}

\begin{theorem}[All original global canonical coefficients are original ring elements]
\label{sc0-proper-canonical-regularity:all-original-sections}
\lean{EnvelopeProperty.Char0.ProperCanonicalCoefficientRegularity.exists_originalRing_coefficient}
\leanok
Let the original regular domain $R$ have its SAME prescribed fraction
field $F$. Let the original integral $Y\to\operatorname{Spec}R$
be universally closed and an isomorphism on an actual nonempty open.
Both original schemes are smooth of the same relative dimension
over the original field, and the original source differentials
have an actual basis. Normalize the original relative canonical
line by its actual nonzero differential section. Then the SAME
prescribed rational coefficient of EVERY original global section
of that entire line belongs to the SAME original $R\subset F$.
The nonzero condition here only permits the original normalization;
the actual BP consumer derives it from its original isomorphism open.
No DVR family, lift, coefficient range or differential-unit
certificate is supplied.
\end{theorem}
\begin{proof}
For EVERY original prime whose actual prime localization is a DVR,
take its literal localization subalgebra inside the SAME $F$.
The true localization algebra equivalence transports the actual
DVR structure to this literal subalgebra, retaining the original
$R$-map and fraction field. Universal closedness gives an actual
lift of its spectrum with BOTH original structural and prescribed
generic triangles. Genuine source localization makes the SAME
entire pulled canonical differential, and then its original
Hom-section with the true unit comparison, isomorphisms.

The preceding all-section local-ring theorem therefore places
the SAME coefficient of EVERY original global section in this
actual original DVR image. Since that image is the literal
localization subalgebra in $F$, these are precisely the membership
conditions in the full original regular-DVR intersection.
The proved original principal-ideal detection theorem identifies
this entire intersection with the SAME original $R$. It produces
an actual $a\in R$ whose image is the original coefficient.
\uses{sc0-valuative-line-regularity:every-original-section,
      sc0-valuative-canonical-section:original-unit,
      sc0-generic-valuative:actual-whole-lifts,
      sc0-regular-dvr-intersection:original-ring}
\end{proof}

This concerns the ENTIRE original relative canonical line's global
sections. The original rounded module is included in that line;
its actual BP consumer and the entire pushforward module comparison
are the next steps. No vanishing conclusion is asserted.
